\documentclass[11pt]{article}
\pagestyle{plain}

% \pgfplotsset{compat=1.18}
\usepackage[margin=1in]{geometry}
\usepackage{amsmath,amsfonts,amssymb,amsthm}
\usepackage[ruled,vlined,linesnumbered,commentsnumbered]{algorithm2e}
\SetKwInput{KwIn}{Input}
\SetKwInput{KwOut}{Output}
\SetKwInput{KwGlobal}{Global variable}
\SetKwInput{KwOracle}{Oracle}
\SetKwComment{Comment}{// }{}
\usepackage{graphicx}
\usepackage{enumerate}
\usepackage{bbm}
\usepackage{verbatim}
\usepackage{hyperref,color}
\usepackage[capitalize,nameinlink]{cleveref}
\usepackage[dvipsnames]{xcolor}
\hypersetup{
	colorlinks=true,
	pdfpagemode=UseNone,
	citecolor=OliveGreen,
	linkcolor=NavyBlue,
	urlcolor=Magenta,
	pdfstartview=FitW
}
\usepackage{appendix}
\usepackage{makecell}
\usepackage{tablefootnote}
\crefname{appsec}{Appendix}{Appendices}
\usepackage{tikz}
\usepackage{pgfplots}
\usepackage{xifthen}
\usepackage{subcaption}
\usepackage{hhline}
\usepackage{listings}

\theoremstyle{plain}
\newtheorem{theorem}{Theorem}[section]
\newtheorem{proposition}[theorem]{Proposition}
\newtheorem{lemma}[theorem]{Lemma}
\newtheorem{corollary}[theorem]{Corollary}
\newtheorem{conjecture}[theorem]{Conjecture}
\newtheorem{problem}[theorem]{Problem}
\newtheorem{claim}[theorem]{Claim}
\newtheorem{fact}[theorem]{Fact}

\theoremstyle{definition}
\newtheorem{definition}[theorem]{Definition}
\newtheorem{example}[theorem]{Example}
\newtheorem{observation}[theorem]{Observation}
\newtheorem{assumption}[theorem]{Assumption}
\newtheorem*{assumption*}{Assumption}
\newtheorem{condition}[theorem]{Condition}

\theoremstyle{remark}
\newtheorem{remark}[theorem]{Remark}

\crefname{lemma}{Lemma}{Lemmas}
\crefname{theorem}{Theorem}{Theorems}
\crefname{definition}{Definition}{Definitions}
\crefname{fact}{Fact}{Facts}
\crefname{claim}{Claim}{Claims}
\crefname{proposition}{Proposition}{Propositions}


\newcommand{\dif}{\,\mathrm{d}}
%\newcommand{\E}{\mathbb{E}}
%\newcommand{\Var}{\mathrm{Var}}
\newcommand{\Cov}{\mathrm{Cov}}
\newcommand{\MI}{\mathrm{I}}
%\newcommand{\Ent}{\mathrm{Ent}}
\newcommand{\ind}{\mathbbm{1}}
\newcommand{\allone}{\mathbf{1}}
\newcommand{\tr}{\mathrm{tr}}
\newcommand{\wt}{\mathrm{weight}}
\DeclareMathOperator*{\argmax}{arg\,max}
\newcommand{\norm}[1]{\left\lVert #1 \right\rVert}
\newcommand{\TV}[2]{\left\|#1 - #2\right\|_{\mathrm{TV}}}
\newcommand{\ceil}[1]{\left\lceil #1 \right\rceil}
\newcommand{\floor}[1]{\left\lfloor #1 \right\rfloor}
\newcommand{\ip}[2]{\left\langle #1 , #2 \right\rangle}
\newcommand{\grad}{\nabla}
\newcommand{\hessian}{\nabla^2}
\newcommand{\trans}{\intercal}
\newcommand{\diag}{\mathrm{diag}}
\newcommand{\poly}{\mathrm{poly}}
\newcommand{\dist}{\mathrm{dist}}
\newcommand{\sgn}{\mathrm{sgn}}
\newcommand{\fpras}{\mathsf{FPRAS}}
\newcommand{\fpaus}{\mathsf{FPAUS}}
\newcommand{\fptas}{\mathsf{FPTAS}}

\newcommand{\eps}{\varepsilon}

\newcommand{\N}{\mathbb{N}}
\newcommand{\Z}{\mathbb{Z}}
\newcommand{\Q}{\mathbb{Q}}
\newcommand{\R}{\mathbb{R}}
\newcommand{\C}{\mathbb{C}}

\renewcommand{\AA}{\mathcal{A}}
\newcommand{\BB}{\mathcal{B}}
\newcommand{\CC}{\mathcal{C}}
\newcommand{\DD}{\mathcal{D}}
\newcommand{\EE}{\mathcal{E}}
\newcommand{\FF}{\mathcal{F}}
\newcommand{\GG}{\mathcal{G}}
\newcommand{\HH}{\mathcal{H}}
\newcommand{\II}{\mathcal{I}}
\newcommand{\JJ}{\mathcal{J}}
\newcommand{\KK}{\mathcal{K}}
\newcommand{\LL}{\mathcal{L}}
\newcommand{\MM}{\mathcal{M}}
\newcommand{\NN}{\mathcal{N}}
\newcommand{\OO}{\mathcal{O}}
\newcommand{\PP}{\mathcal{P}}
\newcommand{\QQ}{\mathcal{Q}}
\newcommand{\RR}{\mathcal{R}}
\renewcommand{\SS}{\mathcal{S}}
\newcommand{\TT}{\mathcal{T}}
\newcommand{\UU}{\mathcal{U}}
\newcommand{\VV}{\mathcal{V}}
\newcommand{\WW}{\mathcal{W}}
\newcommand{\XX}{\mathcal{X}}
\newcommand{\YY}{\mathcal{Y}}
\newcommand{\ZZ}{\mathcal{Z}}

\newcommand{\KL}[2]{D_{\mathrm{KL}} \left( #1 \,\Vert\, #2 \right)}
\newcommand{\KLbig}[2]{D_{\mathrm{KL}} \left( #1 \,\big\Vert\, #2 \right)}
\newcommand{\KLBig}[2]{D_{\mathrm{KL}} \left( #1 \,\Big\Vert\, #2 \right)}

\newcommand{\SKL}[2]{D_{\mathrm{SKL}} \left( #1 \,\Vert\, #2 \right)}
\newcommand{\SKLbig}[2]{D_{\mathrm{SKL}} \left( #1 \,\big\Vert\, #2 \right)}
\newcommand{\SKLBig}[2]{D_{\mathrm{SKL}} \left( #1 \,\Big\Vert\, #2 \right)}

\newcommand{\qandq}{\quad\text{and}\quad}
\newcommand{\supp}{\mathrm{supp}}

\newcommand{\DKL}[2]{\-D_{\-{KL}}\left(#1 \parallel #2\right)}
\newcommand{\DTV}[2]{\-D_{\mathrm{TV}}\left({#1},{#2}\right)}
% \newcommand{\dist}{\mathrm{dist}}
\newcommand{\e}{\mathrm{e}}
\renewcommand{\epsilon}{\varepsilon}
\renewcommand{\emptyset}{\varnothing}
% \newcommand{\norm}[1]{\left\Vert#1\right\Vert}
\newcommand{\set}[1]{\left\{#1\right\}}
\newcommand{\tuple}[1]{\left(#1\right)} 
% \newcommand{\eps}{\varepsilon}
\newcommand{\inner}[2]{\left\langle #1,#2\right\rangle}
\newcommand{\tp}{\tuple}
\newcommand{\ol}{\overline}
\newcommand{\abs}[1]{\left\vert#1\right\vert}
\newcommand{\ctp}[1]{\left\lceil#1\right\rceil}
\newcommand{\ftp}[1]{\left\lfloor#1\right\rfloor}
\newcommand{\Ind}[1]{\-{Ind}_{#1}}
\newcommand{\todo}[1]{{\color{red} TODO: #1}}

\newcommand{\Cor}{\Psi^{\mathrm{sym}}}
\newcommand{\cor}{\Psi^{\-{cor}}}
\newcommand{\Inf}{\Psi}
% \newcommand{\diag}{\-{diag}}

\def\*#1{\boldsymbol{#1}} % Use \*A for \mathbf{A}
\def\+#1{\mathcal{#1}} % Use \+A for \mathcal{A}
\def\-#1{\mathrm{#1}} % Use \-A for \mathrm{A}
\def\=#1{\mathbb{#1}} % Use \^A for \mathbb{A}
\def\!#1{\mathfrak{#1}} % Use \!A for \mathfrak{A}

\newcommand*\diff{\mathop{}\!\mathrm{d}}

% probability
% \DeclareMathOperator*{\oPr}{\mathbf{Pr}}
\def\oPr{\mathop{\mathrm{Pr}}}
\renewcommand{\Pr}[2][]{ \ifthenelse{\isempty{#1}}
  {\oPr\left[#2\right]}
  {\oPr_{#1}\left[#2\right]} } % Use \Pr[a]{b} for \mathbf{Pr}_a[b], \Pr{b} for  \mathbf{Pr}[b]

% expectation: relies on the package "dsfont"
% \DeclareMathOperator*{\oE}{\mathbb{E}}
\def\oE{\mathop{\mathbb{E}}}
\newcommand{\E}[2][]{ \ifthenelse{\isempty{#1}}
  {\oE\left[#2\right]}
  {\oE_{#1}\left[#2\right]} }

% variance
%\DeclareMathOperator*{\oVar}{\mathbf{Var}}
\def\oVar{\mathrm{Var}}
\newcommand{\Var}[2][]{ \ifthenelse{\isempty{#1}}
  {\oVar\left[#2\right]}
  {\oVar_{#1}\left[#2\right]} }

% entropy
% \DeclareMathOperator*{\oEnt}{\mathbf{Ent}}
\def\oEnt{\mathrm{Ent}}
\newcommand{\Ent}[2][]{ \ifthenelse{\isempty{#1}}
  {\oEnt\left[#2\right]}
  {\oEnt_{#1}\left[#2\right]} }

\newcommand{\PhiEnt}[2][]{ \ifthenelse{\isempty{#1}}
  {\oEnt^\phi\left[#2\right]}
  {\oEnt^\phi_{#1}\left[#2\right]} }


\newenvironment{Exampleparagraph}[1][]
{
    \par 
    \noindent 
    \textbf{\color{blue} #1}
    \par 
    \vspace{0.5em} 
    \itshape 
}
{
    \par 
    \vspace{1em} 
}

%\input{local-macro}

\newcommand{\RED}[1]{\textcolor{red}{#1}}

\newcommand{\bern}[1]{\mathrm{Bern}\left(#1\right)}
\newcommand{\pll}{\parallel}



\title{MATH 6112. FALL 2026. ASSIGNMENT 1.}
\date{Due on September 16, 2026}
\author{{\textbf{Zejia Chen}} \\ \textbf{GTID: 904178159}}

\pgfplotsset{compat=1.18} 
\begin{document}

\maketitle
\section*{Problem 1}
\subsection*{Exercise 1.3.13-1}
Suppose there are two different minimal polynomials $p_1,p_2$ of $T$.
We use $\deg f$ to denote the degree of polynomial $f$.
By minimality, $\deg p_1 = \deg p_2$.
Since both $p_1$ and $p_2$ are monic, we have that $\deg (p_1-p_2) < \deg p_1$ and $(p_1-p_2)(T) = p_1(T)-p_2(T) = 0$.
Since $p_1-p_2\neq 0$, it contradicts the minimality of $p_1$.
Therefore, the minimal polynomial is unique.
\subsection*{Exercise 1.3.13-2}
Let $p = a_0 + a_1x + \cdots + a_kx^k$ be the minimal polynomial ($a_k=1$).
It is trivial that $k\geq 1$.
\\$\boldsymbol{\implies}:$ Suppose $a_0 = 0$. 
Then we have that
$$
    0 = T^{-1}p(T) = \sum_{i=1}^k a_i T^{i-1} =: q(T)
$$
where $q(x) = a_1+a_2x+\cdots+a_k x^{k-1}$.
$q\neq 0$ since $k\geq 1$.
$q$ has lower degree than $p$ and $q(T) = 0$.
This contradicts the minimality of $p$.
Hence, $a_0\neq 0$.
\\$\boldsymbol{\impliedby}:$
Since $p(T) = 0$, we have that
$$
    T\tp{-\sum_{i=1}^k \frac{a_i}{a_0}T^{i-1}} = I .
$$
We complete the proof by setting $T^{-1} = -\sum_{i=1}^k \frac{a_i}{a_0}T^{i-1}$.
\subsection*{Exercise 1.3.13-6}
\textbf{(a)}: $Tv = \alpha v$ implies that $0 = T^k v = \alpha T^{k-1}v =\cdots = \alpha^k v$.
Since $v\neq 0$, $\alpha^k = 0$ shows that $\alpha = 0$.
\\\textbf{(b)}: WLOG, we may assume $p=k-1$ ($T^k=0$ and if \(p<k-1\), we set \(\alpha_i=0\) for \(p<i\le k-1\)).
Construct $\{\beta_i\}_{i\in [k-1]\cup\{0\}}$ as follows,
\begin{equation}\label{eq:bP1_beta}
    \beta_0 = \frac 1{\alpha_0}\qquad\text{and}\qquad\forall i\in[k-1], \beta_i = -\frac 1{\alpha_0}\sum_{j=1}^i \alpha_j\beta_{i-j}.
\end{equation}
Set $P = \sum_{i=0}^{k-1}\beta_i T^i$.
By \Cref{eq:bP1_beta}, $SP = I + T^kH = I$, where $H$ is some polynomial of $T$.
Therefore, $S$ is invertible and $S^{-1}=P$.
\section*{Problem 2}
We first prove that $C$ is invertible.
Let $V = (v_1,\cdots,v_n)$ and $W = (w_1,\cdots,w_n)$.
Let $D$ be the matrix associated to the linear transformation which changes the basis: $v_j = \sum_{i=1}^n d_{ij}w_i$.
We have that $VC = W$ and $WD = V$.
Therefore, $VCD = WD = V$.
Since $v_1\ldots v_n$ are linearly independent, $CD = I$ and $C$ is invertible.

We then prove that $B=DAC$.
The following holds:
\begin{align*}
    WDAC &= VAC = [T(v_1),\cdots T(v_n)]C &&(\text{def. of }A)
    \\&= [\sum_{i=1}^n c_{i1} T(v_i), \sum_{i=1}^n c_{i2} T(v_i) ,\cdots , \sum_{i=1}^n c_{in} T(v_i)]
    \\&= [T(\sum_{i=1}^n c_{i1} v_i), T(\sum_{i=1}^n c_{i2}v_i) ,\cdots , T(\sum_{i=1}^n c_{in}v_i)] &&(T\in \-{Hom}(V,V))
    \\&= [T(w_1),\cdots T(w_n)] = WB.  &&(\text{def. of }B)
\end{align*}
Since $w_1,\cdots,w_n$ are linearly independent, we have that $DAC = B$.
\section*{Problem 3}
\textbf{(a)}:
By Block Diagonal Form (Corollary 1.5.45), there exists $V$ such that $V^{-1}AV = R = \diag\{R_{ii}\}_{i\in [p]}$, where $R_{ii} = \lambda_i I + U_{ii}$ and $U_{ii}\in \C^{m_i\times m_i}$ is strictly upper triangular.
Fix $k\in [p]$, let $U_{kk} = (u_{ij})_{i,j\in [m_k]}$ and $v_1\ldots v_{m_k}$ be the $(1+\sum_{i<k}m_i)$-th to $(\sum_{i\leq k}m_i)$-th column vectors of $V$.
Since $AV = VR$, we have that
$$
    \begin{cases}
        Av_1 = \lambda_k v_1
        \\Av_2 = u_{12}v_1 + \lambda_k v_2
        \\\ \ \ \ \ \vdots
        \\Av_{m_k} = \lambda_k v_{m_k} + \sum_{i<m_k} u_{i m_k}v_i
    \end{cases}
    \implies 
    \begin{cases}
        (A - \lambda_k I)v_1 = 0
        \\ (A - \lambda_k I)^2 v_2 = u_{12} (A - \lambda_k I)v_1 = 0
        \\\ \ \ \ \ \vdots
        \\ (A - \lambda_k I)^{m_k}v_{m_k} = (A - \lambda_k I)^{m_k-1}\sum_{i<m_k} u_{i m_k}v_i = 0
    \end{cases}
$$
Therefore, $v_1,\cdots,v_{m_k}$ are solutions of $(A-\lambda_k I)^{m_k}v = 0$ and the solution set contains at least $m_k$ linearly independent vectors.
On the other hand, let $v$ be one column vector of $V$ and $v\notin \{v_1,\cdots,v_{m_k}\}$.
Suppose $\lambda$ be the corresponding eigenvalue of $v$, we have that 
$$
    (A-\lambda_k I)^{m_k}V=V(R - \lambda_k I)^{m_k} \implies
$$
$$
    (A-\lambda_k I)^{m_k}v = (\lambda - \lambda_k)^{m_k}v + \text{linear combination of column vectors of $V$ except $v$} \neq 0,
$$ 
since $\lambda - \lambda_k\neq 0$ and column vectors of $V$ are linearly independent.
Since $\-{rank}(V)=n$, the dimension of solution set of $(A-\lambda_k I)^{m_k}v = 0$ is $m_k$.
\\\textbf{(b)}
Putting together the sets of vectors from \textbf{(a)}, we get exactly $V$. 
Since $\-{rank}(V)=n$, it is a basis of $\C^n$.
\section*{Problem 4}
\textbf{(a)}
Suppose $v_1,\cdots v_n$ are column vectors of $V$.
$A V =V J_n(\lambda)$ implies that 
$$
    \forall i\in[n], (A - \lambda I)v_i = v_{i-1} \quad \text{where} \quad v_0 = 0.
$$
Since $(A - \lambda I)^{n-1} = V (J_n(\lambda) - \lambda I)^{n-1}V^{-1}\neq 0$, it is trivial to find a vector $x$ such that $(A - \lambda I)^{n-1}x \neq 0$.
Set $v_i = (A-\lambda I)^{n-i}x$.
It is sufficient to show the linear independence of $\{v_i\}_{i\in[n]}$.
Suppose not, that is, there exists $j \leq n$ such that
$$
    v_j = \sum_{i=1}^{j-1} a_i v_i \implies (A-\lambda I)^{n-j}x = \sum_{i=1}^{j-1}a_i (A-\lambda I)^{n-i}x
    \implies (A-\lambda I)^{n-1}x = \sum_{i=1}^{j-1}a_i (A-\lambda I)^{n-i+j-1}x.
$$
Since $(A - \lambda I)^{n} = V (J_n(\lambda) - \lambda I)^nV^{-1} = 0$, we have that
$$
    (A - \lambda I)^{n-1}x = (A-\lambda I)^n\sum_{i=1}^{j-1}a_i (A-\lambda I)^{j-i-1}x = 0,
$$
which contradicts our choice of $x$. 
Therefore, $\{v_i\}_{i\in[n]}$ are linearly independent and $V = [v_1,\cdots v_n]$.
\\\textbf{(b)}
Suppose there is only one eigenvalue $\lambda$ and the Jordan blocks have sizes $m_1 \geq \cdots \geq m_k$.
We get the corresponding vectors for $V$ as follows:
\begin{itemize}
    \item [1.] Initially $i=1$ and $S = \emptyset$.
    \item [2.] Take $x_i\in \ker (A-\lambda I)^{m_i}\setminus \ker (A-\lambda I)^{m_i-1}$ and $(A - \lambda I)^{m_i-1}x_i \notin \-{span}(S)$.
    Let the vectors corresponding to the $i$-th Jordan block to be 
    $$
        V_i =  [(A - \lambda I)^{m_i-1}x_i, (A - \lambda I)^{m_i-2}x_i,\cdots,(A - \lambda I)x_i,x_i].
    $$
    \item [3.] Set $S \leftarrow S \cup \{(A - \lambda I)^{m_i-1}x_i\}$, $i\leftarrow i+1$ and back to Step 2 until $i > k$.
    \item [4.] $V = [V_1,\ldots,V_k]$.
\end{itemize}
Note that such $x$ always exists since $\dim \ker(A-\lambda I)^{m_i} - \dim \ker(A-\lambda I)^{m_i-1} \geq i$ (there are at least $i$ Jordan blocks whose sizes are at least $m_i$) and $\dim\-{span} (S) < i$ at step $i$.
Finding $x_i$ can be achieved by Gaussian elimination.


By \textbf{(a)}, the choice of $V$ satisfies $AV = V \diag\{J_{m_i}(\lambda)\}_{i\in[k]}$ and for $i\in [k]$, the columns of $V_i$ are linearly independent. 
Suppose the columns of $V$ are not linearly independent, that is, WLOG, we have that there exists $\*a\neq \*0$ (we can get rid of $(A-\lambda I)^{\ell}x_i$ for $\ell < m_i-1$ by multiplying $(A-\lambda I)^{m_i-1-\ell}$ on both sides),
$$
    \sum_{i=1}^k a_i (A-\lambda I)^{m_i - 1} x_i = 0,
$$
which contradicts our Step 2.
In conclusion, $V$ is a feasible choice.
Also, note that for each column $v$ of $V$, $v$ is the solution of $(A-\lambda I)^{m_1+\cdots+m_k} u = 0$.

For Jordan blocks with different eigenvalues, by Problem 3, we complete the proof by putting all $V$'s of the different eigenvalues together.
\section*{Problem 5}
We prove by induction on $k$. 
The statement follows from definition of the Jordan block if $k=1$.

Now suppose $k > 1$.
We have that 
\begin{align}
    (J^k)_{ij} &= \sum_{\ell=1}^n  (J^{k-1})_{i\ell} (J)_{\ell j}
    =  \sum_{\ell=i}^{\min(n,i+k-1)} \binom{k-1}{\ell - i}\lambda^{k-1-(\ell-i)} (J)_{\ell j} &&(\text{I.H.}) \nonumber
    \\&= \lambda^{k+i-j} \sum_{\ell=i}^{\min(n,i+k-1)} \ind[j-1\leq \ell \leq j]\binom{k-1}{\ell - i} &&(\text{def. of }J) \nonumber
    \\&= \lambda^{k+i-j}\times 
    \begin{cases}
        \binom{k-1}{j-i-1} + \binom{k-1}{j-i}&,i\leq j-1 \land j\leq i+k-1
        \\\binom{k-1}{j-i}=1=\binom{k}{j-i}&, j=i
        \\\binom{k-1}{j-i-1}=1=\binom{k}{j-i}&, j=i+k
        \\0 &, \text{o.w.}
    \end{cases} . \label{eq:P5_final}
\end{align}
Since $\binom{k-1}{j-i-1} + \binom{k-1}{j-i} = \binom{k}{j-i}$, \Cref{eq:P5_final} completes the proof.
\section*{Problem 6}
By definition we have that
\begin{align*}
    \frac{d(\det A(\mu))}{d\mu}
    &= \sum_{\sigma \in S} (-1)^\sigma \frac{d}{d\mu} \prod_{i=1}^n a_{\sigma_i,i}
    = \sum_{\sigma \in S} (-1)^\sigma \sum_{j=1}^n a_{\sigma_j,j}' \prod_{i\in[n],i\neq j}^n a_{\sigma_i,i}
    \\&= \sum_{j=1}^n \sum_{\sigma \in S} (-1)^\sigma a_{\sigma_j,j}' \prod_{i\in[n],i\neq j}^n a_{\sigma_i,i}
    = \sum_{j=1}^n \det (a_1,\cdots,a_{j-1},a_j',a_{j+1},\cdots,a_n).
\end{align*}
\section*{Problem 7}
Block Diagonal Form (Corollary 1.5.45) shows the existence of $T$.
We now prove the uniqueness property.
Suppose $T^{-1}AT = \diag\{B_{ii}\}_{i\in[n]}$ and $(T')^{-1}AT' = \diag\{B_{ii}'\}_{i\in[p]}$ where $B_{ii}$ has the same eigenvalues as $B_{ii}'$ for $i\in[p]$.
It is sufficient to show that $T = T'V$ and $V$ is an diagonal block matrix.
Note that the invertiblity of $V$ follows from the invertiblity of $T$ and $T'$. 

Let $C = (T')^{-1}T$. 
We write $C = (C_{ij})_{i,j\in[p]}$ as a block matrix.
We need to show that $C_{ij} = 0$ for $i\neq j$.
We have that
\begin{equation}\label{eq:P7_1}
    T \diag\{B_{ii}\}_{i\in [p]}T^{-1} = (T') \diag\{B_{ii}'\}_{i\in [p]}(T')^{-1} \implies C\cdot \diag\{B_{ii}\}_{i\in [p]} = \diag\{B_{ii}'\}_{i\in [p]}\cdot C .
\end{equation}
Therefore, \Cref{eq:P7_1} shows that
$$
    \begin{pmatrix}
        C_{11}B_{11}   &\cdots  &C_{1p}B_{pp}
        \\\vdots &\ddots  &\vdots
        \\C_{p1}B_{11} &\cdots  &C_{pp}B_{pp}
    \end{pmatrix}
    = \begin{pmatrix}
        B_{11}'C_{11}  &\cdots  &B_{11}'C_{1p}
        \\\vdots &\ddots  &\vdots
        \\B_{pp}'C_{p1} &\cdots  &B_{pp}'C_{pp}
    \end{pmatrix} .
$$
Fix $i,j\in[p]$ and $i\neq j$. 
Suppose $V^{-1} B_{jj} V =  R_{jj}$ and $W^{-1} B'_{ii} W = R'_{ii}$ where $R_{jj}$ and $R'_{ii}$ are upper-triangular matrix and lower-triangle matrix respectively.
Let $H = W^{-1}C_{ij} V$, we have that 
$$
    C_{ij} B_{jj} = B_{ii}'C_{ij} \implies W^{-1}C_{ij}VV^{-1}B_{jj}V = W^{-1}B_{ii}'WW^{-1}C_{ij}V \implies HR_{jj} = R_{ii}'H.
$$
Let $\lambda_1 \cdots \lambda_{n_j}$ be diagonal entries of $R_{jj}$ and $\mu_1 \cdots \mu_{n_i}$ be diagonal entries of $R_{ii}'$.
Since the eigenvalues of $R_{jj}$ and $R_{ii}'$ are disjoint, $\{\lambda_k\}_{k\in[n_j]} \cap \{\mu_k\}_{k\in[n_i]} = \emptyset$. 
$HR_{jj} = R_{ii}' H$ implies the following equations
$$
    \lambda_1 h_{11} = \mu_1 h_{11}\implies h_{11}=0,\quad h_{11}r_{12}+h_{12}\lambda_2 = \mu_1h_{12}\implies h_{12}=0\cdots .
$$
By the the argument in the class, $H = 0$.
Since $V,W$ are of full rank, $H = 0$ implies $C_{ij} = 0$.
\end{document}